% !TeX spellcheck = it_IT

\documentclass[format=169, handout, 10pt]{beamer}

\usepackage[utf8]{inputenc}
\usepackage[italian]{babel}
\usepackage{hyperref}
\usetheme{metropolis}

\usepackage{mathtools}
\DeclarePairedDelimiter{\ceil}{\lceil}{\rceil}

\title{Esercitazioni di Informatica B}
\subtitle{Struct in MATLAB}

\author{Stefano Cereda\\
	stefano.cereda@polimi.it
}
\date{27/11/2018}
\institute[PoliMi]{\vspace{0.5cm}\centering Politecnico di Milano \\ \vspace{0.2cm}
	\includegraphics[width=0.2\textwidth]{../logopolimi}}

\setbeamercovered{invisible}


\makeindex


\begin{document}
\begin{frame}
	\maketitle
\end{frame}


\begin{frame}{Paghe}
Data la seguente tabella:

\begin{tabular}{|c|c|c|c|c|c|}
	\hline
	& op1 & op2 & op3 & op4 & op5 \\
	\hline
	Paga oraria & 5 & 5.5 & 6.5 & 5 & 6.25 \\ 
	\hline
	Ore settimanali & 40 & 43 & 37 & 50 & 45 \\
	\hline
	Pezzi prodotti & 1000 & 1100 & 1000 & 1200 & 1100 \\
	\hline
\end{tabular} 

Scrivere uno script Matlab che organizzi la tabella in un array di dati strutturati e risponda alle seguenti domande:
\begin{enumerate}
	\item Quanto guadagna ogni operaio?
	\item Qual è il salario totale di tutti gli impiegati?
	\item Quanti pezzi vengono prodotti?
	\item Qual è il costo medio di un pezzo?
	\item Quante ore occorrono in media per un pezzo?
	\item Qual è l’operaio più efficiente?
\end{enumerate}
\end{frame}

\begin{frame}[allowframebreaks]{Archivio mammiferi}

	Si vogliono rappresentare informazioni relative ai mammiferi ospitati in un
	parco protetto.
	
	In particolare, per ogni mammifero si rappresentano i seguenti dati:\\
	NOME TIPO, TIPO ALIMENTAZIONE, PESO MEDIO, insieme di ESEMPLARI
	
	Dove ogni esemplare è caratterizzato da:\\
	CODICE ESEMPLARE, ANNO DI NASCITA, NUMERO CUCCIOLI
	
	\begin{enumerate}
	\item Si crei una variabile mammifero che contenga le informazioni del
	seguente mammifero:\\
	NOME TIPO: Elefante\\
	TIPO ALIMENTAZIONE: Erbivoro\\
	PESO MEDIO: 2500\\
	ESEMPLARE 1: (Codice: 3, Anno: 1985, Cuccioli: 3)\\
	ESEMPLARE 2: (Codice: 5, Anno: 1989, Cuccioli: 0)\\
	ESEMPLARE 3: (Codice: 8, Anno: 1982, Cuccioli: 0)\\
	
	\item Chieda all'utente se vuole inserire ulteriori elefanti
	\item Si scriva il codice che costruisce, a partire dalla variabile creata al punto
	precedente, l'archivio degli esemplari di elefante sterili sapendo che un
	elefante si dice sterile se è nato prima del 1990 e non ha concepito
	nessun cucciolo.
\end{enumerate}

	Si stampi su schermo l'elenco dei codici degli elefanti sterili.
	
\end{frame}

\iffalse
\begin{frame}[allowframebreaks]{Treni in stazione}
Si consideri il sistema di gestione dei treni in una stazione. In ogni stazione ci sono varie banchine e
ogni banchina consente la sosta e il transito di un treno. Ad ogni treno in stazione viene assegnato uno dei
seguenti stati:
\begin{itemize}
\item fuoriStazione: il treno, in attesa di essere assegnato a una banchina, si trova fuori dalla stazione
\item inIngresso: il treno è stato assegnato a una banchina ed è in marcia a velocità ridotta per entrare in
stazione
\item inSosta: il treno è fermo a una banchina
\item attesaOut: il treno è fermo a una banchina in attesa di poter partire
\item  inUscita: il treno sta abbandonando la banchina procedendo a velocità ridotta.
\end{itemize}
Si assuma che siano già stati introdotti i tipi Stazione e Banchina. In particolare, Stazione contiene l'insieme delle banchine presenti in stazione (banchine), la coda dei treni in attesa all'ingresso (codaTreni) e una variabile Booleana che dice se c'è qualche treno in manovra in stazione (bloccata).

Banchina, invece, contiene il numero della banchina e i dati del
treno in sosta, se la banchina è in stato occupato. Ciascun treno, oltre ad avere il proprio stato e il numero della banchina eventualmente assegnata, ha un campo minutiAttesaOut che indica da quanti minuti è in
attesa in stato attesaOut.

Si supponga che sia stata definita la variabile stazioneMI di tipo Stazione e che questa variabile sia stata
riempita in modo opportuno per rappresentare lo stato corrente della stazione di Milano, si definisca uno script che, se la stazione non è bloccata dalle manovre di qualche
treno e se vi è almeno un treno in coda per entrare in stazione, blocca la stazione, estrae dalla coda il
primo treno che vi era stato inserito, gli assegna lo stato inIngresso, e aggiorna la coda in modo tale
che poi contenga solo i rimanenti treni in coda. Se per esempio, prima di questa operazione la coda è cosi costituita: \\
stazioneMI.codaTreni(1) = FR2092, stazioneMI.codaTreni(2) = IR097,
stazioneMI.codaTreni(3) = FA2673 (per semplicità abbiamo riportato solo il nome del treno, ma,
ovviamente, codaTreni dovrebbe contenere l’intera struct corrispondente), dopo l’operazione si deve
avere che stazioneMI.codaTreni(1)= IR097, stazioneMI.codaTreni(2) = FA2673.
\end{frame}
\fi

\end{document}